\section{The Propers: Notable and Quotable}

Few pericopes in the temporal cycle have had a stranger second career than this
one. The gallery below collects five documented later uses of wording from the
appointed Gospel in which the register or force of the phrase is genuinely
turned --- into satire, into anti-clerical polemic, into a swindler's private
code, into a theory of political language, into a parliamentary insult.
Straight exegesis and devotional citation are excluded; those belong to the
commentary above.

One philological note governs the whole gallery. The English phrases that
entered the language are the Authorised Version's --- ``the mammon of
unrighteousness,'' ``everlasting habitations,'' ``wiser in their generation''
--- not the Douay's ``mammon of iniquity'' and ``everlasting dwellings,'' and
not the Latin of the missal. The afterlives below therefore attach to a
Protestant rendering of the verses this Mass appoints. Where Latin is given it
is the appointed text; where an English phrase is quoted it is the later
author's own wording, and the difference is part of the evidence.

\begin{afterlife}{1. The swindler's only regret \normalfont(\emph{Gosp.}, Lk 16:9)}
\afield{Appointed phrase}{\latin{fácite vobis amícos de mammóna iniquitátis}}
\afield{Later use}{Anthony Trollope, \work{The Way We Live Now} (1875), ch.\
LXXXI, ``Mr.\ Cohenlupe Leaves London.''}
\afield{Exact wording}{``But why had he, so unrighteous himself, not made
friends to himself of the Mammon of unrighteousness? Why had he not conciliated
Lord Mayors? \dots\ Why had he been insolent at the India Office?''}
\afield{Context}{Augustus Melmotte, the collapsing financier, walks toward
Westminster on the afternoon before his suicide, conducting a self-examination.
Trollope has just told the reader that ``never once, not for a moment, did it
occur to him that he should repent of the fraud in which his whole life had
been passed.''}
\afield{The turn}{Christ's counsel becomes the swindler's single article of
contrition --- and a technical one. Melmotte does not regret the frauds; he
regrets a failure of network management. The verse Augustine laboured to keep
from licensing plunder is here quoted, unsignalled and accurately, by a man
whose only moral standard is that he did not cultivate enough useful friends.
The list that follows --- Lord Mayors, the India Office, Abchurch Lane ---
supplies the parable's ``debtors'' with a Victorian address book.}
\afield{Dependence and rights}{Echo: exact Authorised Version wording,
capitalised, unattributed. Public domain; read in the Project Gutenberg text
(ebook 5231), 2026-07-25.}
\end{afterlife}

\begin{afterlife}{2. ``And he was commended.'' The House of Commons, 1931
\normalfont(\emph{Gosp.}, Lk 16:8--9)}
\afield{Appointed phrases}{\latin{laudávit dóminus víllicum iniquitátis};
\latin{recípiant vos in ætérna tabernácula}}
\afield{Later use}{Sir Hilton Young MP, House of Commons, Representation of the
People (No.\ 2) Bill, 3 February 1931, HC Deb vol.\ 247, cc.\ 1662--1663, with
an interjection by Mr Ernest Brown; the Prime Minister, Ramsay MacDonald,
replies at c.\ 1664.}
\afield{Exact wording}{Young: ``Let me remind hon.\ Members opposite and the
Prime Minister of the parable of the unjust steward; how he had wasted his
master's substance; how he was to be called to account; and how he busied
himself in making friends with the mammon of unrighteousness in order that he
might be received into the everlasting habitations---'' \textsc{Mr Ernest
Brown}: ``And he was commended.'' \textsc{Sir H.\ Young}: ``He was commended in
those days, and there are still some who would commend him. But there are no
everlasting habitations for a Government which is prepared to govern without
the will of the people.'' MacDonald, in reply, accused his critics of holding
that reformers are ``guilty of all the evils of the unjust steward and corrupt
politicians.''}
\afield{The turn}{Two turns at once. First, the eschatological
\latin{tabernácula} are demoted to a ministry's tenure of office. Second, and
more remarkable, a backbencher's four-word interjection puts the pericope's
central exegetical crux on the floor of the House --- \emph{he was commended} ---
and Young's improvised answer, ``He was commended in those days, and there are
still some who would commend him,'' is in substance Augustine's \latin{e
contrario} reading delivered as a debating point. The Prime Minister then names
the parable back across the despatch box.}
\afield{Dependence and rights}{Documented on both sides: the parable is named
three times in the debate. Parliamentary copyright; quoted under the Open
Parliament Licence from the Historic Hansard text, read 2026-07-25.}
\end{afterlife}

\begin{afterlife}{3. Kingsley turns two verses into anti-clerical ammunition
\normalfont(\emph{Gosp.}, Lk 16:3 and 16:6--7)}
\afield{Appointed phrases}{\latin{fódere non váleo, mendicáre erubésco};
\latin{accipe cautiónem tuam \dots\ scribe quinquagínta}}
\afield{Later use}{Charles Kingsley, \work{Westward Ho!} (1855), ch.\ XVIII,
``How They Took the Pearls at Margarita.''}
\afield{Exact wording}{Amyas Leigh: ``No wonder that young men \dots\ will not
listen to the Gospel, while it is preached to them by men on whom they cannot
but look down; a set of softhanded fellows who cannot dig, and are ashamed to
beg.'' His brother Frank, immediately after: ``\dots\ Popish priests and friars,
who are vowed not to be men, and get their bread shamefully and rascally by
telling sinners who owe a hundred measures to sit down quickly and take their
bill and write fifty.''}
\afield{The turn}{Two different verses of the same appointed Gospel, redirected
in consecutive speeches by two brothers, against two different clerical castes.
The steward's excuse becomes a muscular-Christian sneer at unmanly Protestant
clergy, his physical incapacity read as effeminacy. Then the steward's
fraudulent write-down becomes a figure for sacramental absolution: the priest
as a crooked accountant discounting sinners' bills. A parable the Fathers used
to defend almsgiving is here made to attack the remission of sins.}
\afield{Dependence and rights}{Documented: ``a hundred measures \dots\ write
fifty'' is unmistakably signalled. Public domain; read in the Project Gutenberg
text (ebook 1860), 2026-07-25.}
\end{afterlife}

\begin{afterlife}{4. A steward at a garden party fears he is \emph{the} steward
\normalfont(\emph{Gosp.}, Lk 16:1--9)}
\afield{Appointed phrases}{\latin{víllicum iniquitátis}; \latin{fácite vobis
amícos de mammóna iniquitátis}}
\afield{Later use}{Anthony Trollope, \work{Barchester Towers} (1857), ch.\
XXXIX, ``The Lookalofts and the Greenacres.''}
\afield{Exact wording}{``But Mr.\ Plomacy was not quite happy in his mind, for
he thought of the unjust steward and began to reflect whether he had not made
for himself friends of the mammon of unrighteousness.''}
\afield{Context}{Mr Plomacy is Miss Thorne's actual steward, policing the
Ullathorne Sports. Wheedled by Farmer Greenacre into letting an uninvited
plasterer stay to the feast, he suffers an instant scruple of conscience out of
Luke 16 --- and then, unruffled, rises to say grace.}
\afield{The turn}{The joke depends on an exact congruence of office and a total
incongruity of scale. Plomacy is a steward; he has just shown irregular leniency
with his employer's provisions to a man who may be useful; the verse fits him
with uncomfortable precision --- and the stake is who gets fed at a village
f\^ete. Trollope thereby stages, comically, the difficulty Ambrose and Augustine
argued over: whether a small act of unauthorised generosity with someone else's
goods is commendable foresight or plain fraud. The novel declines to decide and
serves dinner.}
\afield{Dependence and rights}{Documented: Trollope names the unjust steward.
Public domain; read in the Project Gutenberg text (ebook 3409), 2026-07-25.}
\end{afterlife}

\begin{afterlife}{5. Words as unjust stewards \normalfont(\emph{Gosp.}, Lk 16:1--2)}
\afield{Appointed phrases}{\latin{víllicum}; \latin{redde ratiónem
villicatiónis tu\ae}}
\afield{Later use}{John Ruskin, \work{Sesame and Lilies} (1865), lecture I, ``Of
Kings' Treasuries.''}
\afield{Exact wording}{``There never were creatures of prey so mischievous,
never diplomatists so cunning, never poisoners so deadly, as these masked words;
they are the unjust stewards of all men's ideas: whatever fancy or favourite
instinct a man most cherishes, he gives to his favourite masked word to take
care of for him; the word at last comes to have an infinite power over him.''}
\afield{The turn}{The parable migrates out of economics into what would later be
called ideology-critique. Ruskin's ``masked words'' --- abstractions everyone
uses and no one defines --- are entrusted with meanings, embezzle them, and end
by ruling the man who deposited them. The transfer is exact on every term: the
goods are ideas, the owner is the speaker, the wasting is semantic drift, and
the accounting Ruskin demands of his readers is \latin{redde ratiónem
villicatiónis tu\ae} applied to vocabulary. It is the only use in this gallery
in which the steward is not a person at all.}
\afield{Dependence and rights}{Documented by the phrase ``unjust stewards.''
Public domain; read in the Project Gutenberg text (ebook 1293), 2026-07-25.}
\end{afterlife}

\clearpage
